\subsection{Regular representations}

Let \(p\) be a real number \(\geqslant 1\) and \(\mu'\) a right Haar measure on \(G\) . For every \(g \in G\) and every function \(f \in \mathcal{L}^p(G, \mu)\) , we denote by \(\gamma_G^{(p)}(g)f\) the function \(x \mapsto f(g^{-1}x)\) on \(G\) . The map \(g \mapsto \gamma_G^{(p)}(g)\) is a continuous linear representation of \(G\) in \(\mathcal{L}^p(G)\) . By passing to quotients, it induces a continuous isometric linear representation of \(G\) in \(L^p(G)\) (INT, VIII, p. 135, § 2, no. 5, prop. 8), denoted \(\gamma_G^{(p)}\) .

Similarly, for every \(g \in G\) and every function \(f \in \mathcal{L}^{p}(G, \mu')\) , we denote by \(\boldsymbol{\delta}_{\mathrm{G}}^{(p)}(g)f\) the function \(x \mapsto f(xg)\) on G. The map \(g \mapsto \boldsymbol{\delta}_{\mathrm{G}}^{(p)}(g)\) is a continuous linear representation of G in \(\mathcal{L}^{p}(G, \mu')\) . By passing to quotients, it induces a continuous isometric linear representation of G in \(\mathrm{L}^{p}(\mathrm{G}, \mu')\) (cf. INT, VIII, p. 136, § 2, no. 5).

We say that \(\boldsymbol{\gamma}_{\mathrm{G}}^{(p)}\) is the left regular representation of G in \(\mathcal{L}^{p}(\mathrm{G})\) or \(\mathrm{L}^{p}(\mathrm{G})\) and that \(\boldsymbol{\delta}_{\mathrm{G}}^{(p)}\) is the right regular representation of G in \(\mathcal{L}^{p}(\mathrm{G},\mu^{\prime})\) or \(\mathrm{L}^{p}(\mathrm{G},\mu^{\prime})\) .

Lemma 3. — Let p be a real number ≥ 1. The left (resp. right) regular representation of G in L \(^{p}\) (G,μ) (resp. in L \(^{p}\) (G,μ')) is faithful.

More precisely, let q be the conjugate exponent of p and let g be an element of G such that \(g \neq e\) .

\begin{enumerate}
\def\labelenumi{\alph{enumi})}
\item
  There exists a function \(\varphi\in\mathcal{K}(\mathrm{G})\) , positive and nonzero in \(\mathrm{L}^{q}(\mathrm{G},\mu)\) , such that \(\langle\varphi,\boldsymbol{\gamma}_{\mathrm{G}}^{(p)}(g)\varphi\rangle=0\) ;
\item
  There exists a function \(\varphi\in\mathcal{K}(\mathrm{G})\) , positive and nonzero in \(\mathrm{L}^{q}(\mathrm{G},\mu')\) , such that \(\langle\varphi,\boldsymbol{\delta}_{\mathrm{G}}^{(p)}(g)\varphi\rangle=0\) .
\end{enumerate}

Consider the case of the left regular representation \(\boldsymbol{\gamma}_{\mathrm{G}}^{(p)}\) , that of the right regular representation being similar. The assertion \(a\) ) implies that \(\boldsymbol{\gamma}_{\mathrm{G}}^{(p)}\) is faithful, for if \(g \neq e\) is an element of G, and if \(\varphi\) is as in \(a\) ), one has \(\varphi \neq \boldsymbol{\gamma}_{\mathrm{G}}^{(p)}(g)\varphi\) since \(\langle \varphi, \varphi \rangle = \int_{\mathrm{G}} \varphi^2 > 0\) .

Let us therefore prove \(a\) ). Let \(g \neq e\) in G. Let C be a compact symmetric neighborhood of \(e\) such that \(g \notin \mathrm{C}^2\) and let \(\varphi \in \mathcal{K}(\mathrm{G})\) be a continuous positive function with integral 1 and support contained in C; the function \(\varphi\) is nonzero in \(\mathrm{L}^p(\mathrm{G}, \mu)\) . Since \(\mathrm{C} \cap g^{-1}\mathrm{C} = \emptyset\) , we have

\[
\langle \varphi , \boldsymbol {\gamma} _ {\mathrm{G}} ^ {(p)} (g) \varphi \rangle = \int_ {\mathrm{G}} \varphi (x) \varphi (g ^ {- 1} x) d \mu (x) = 0.
\]

Let \(\varrho\) be a continuous and bounded representation of G in a Banach space E. For every \(f\in \mathrm{L}^1 (\mathrm{G})\) (resp. \(f^{\prime}\in \mathrm{L}^{1}(\mathrm{G},\mu^{\prime}))\) and every \(g\in \mathrm{G}\) , we have

\[
\varrho (g) \varrho (f \cdot \mu) = \varrho (\varepsilon_ {g} * (f \cdot \mu)) = \varrho (\boldsymbol {\gamma} _ {\mathrm{G}} ^ {(1)} (g) f \cdot \mu),\tag{1}
\]

\[
\varrho (f ^ {\prime} \cdot \mu^ {\prime}) \varrho (g) = \varrho ((f ^ {\prime} \cdot \mu^ {\prime}) * \varepsilon_ {g}) = \varrho (\boldsymbol {\delta} _ {\mathrm{G}} ^ {(1)} (g ^ {- 1}) f ^ {\prime} \cdot \mu ^ {\prime}),\tag{2}
\]

(INT, VIII, p. 144, § 3, no. 2, formula (5)).

Suppose G is unimodular, and set \(\mu' = \mu\) . The biregular representation of G in \(\mathcal{L}^{p}(\mathrm{G})\) (resp. in \(\mathrm{L}^{p}(\mathrm{G})\) ) is the representation \(\varrho_{\mathrm{G}}^{(p)}\) of \(G \times G\) in \(\mathcal{L}^{p}(\mathrm{G})\) (resp. in \(\mathrm{L}^{p}(\mathrm{G})\) ) such that

\[
\boldsymbol {\varrho} _ {\mathrm{G}} ^ {(p)} (g _ {1}, g _ {2}) = \boldsymbol {\gamma} _ {\mathrm{G}} ^ {(p)} (g _ {1}) \boldsymbol {\delta} _ {\mathrm{G}} ^ {(p)} (g _ {2}) = \boldsymbol {\delta} _ {\mathrm{G}} ^ {(p)} (g _ {2}) \boldsymbol {\gamma} _ {\mathrm{G}} ^ {(p)} (g _ {1}).
\]

It is a continuous linear representation (lemma 1 of V, p. 377). It satisfies

\[
(\boldsymbol {\varrho} _ {\mathrm{G}} ^ {(p)} (g _ {1}, g _ {2}) f) (x) = f (g _ {1} ^ {- 1} x g _ {2})
\]

for all \(f \in \mathcal{L}^p(\mathrm{G})\) , all \((g_1, g_2) \in \mathrm{G} \times \mathrm{G}\) and every \(x \in \mathrm{G}\) .

Note. --- The biregular representation of G in \(\mathrm{L}^{p}(\mathrm{G},\mu)\) is not necessarily faithful; its kernel is the image of the center of G under the map \(g\mapsto(g,g)\) (exercise 4 of V, p. 487).

When \(p = 2\) the left regular representation \(\gamma_{\mathrm{G}}^{(2)}\) of G in the complex Hilbert space \(\mathrm{L}^2 (\mathrm{G},\mu)\) is unitary, since it is isometric. Likewise, the right regular representation \(\delta_{\mathrm{G}}^{(2)}\) in \(\mathrm{L}^2 (\mathrm{G},\mu')\) is unitary.

We will simply denote \(\gamma_{\mathrm{G}} = \gamma_{\mathrm{G}}^{(2)}\) and \(\delta_{\mathrm{G}} = \delta_{\mathrm{G}}^{(2)}\) and these representations will be called the left and right regular representations of G.

If G is unimodular, the biregular representation \(\varrho_{\mathrm{G}}^{(2)}\) of \(G \times G\) in \(\mathrm{L}^{2}(\mathrm{G}, \mu)\) is unitary. We will simply denote it by \(\varrho_{G}\) .

Lemma 4. --- Let \(p\) be a real number \(\geqslant 1\) . For \(f \in \mathcal{L}^1(\mathrm{G})\) , the linear map \(\gamma_{\mathrm{G}}^{(p)}(f)\) coincides with the endomorphism of \(\mathrm{L}^p(\mathrm{G})\) defined by \(\varphi \mapsto f *^\mu \varphi\) .

This follows from INT, VIII, p. 157, § 4, n°2, prop. 6, taking into account formula (14) of INT, VIII, p. 165.

Note. --- Recall that, for any function \(f\) on G, one defines the function \(\check{f}\) on G by \(\check{f}(g) = f(g^{-1})\) for all \(g \in \mathrm{G}\) (INT, VII, p. 12, § 1, n° 1, formula (12)).

One checks that if \(f \in \mathcal{L}^{1}(\mathrm{G})\) , the linear map \(u = \boldsymbol{\delta}_{\mathrm{G}}^{(p)}(f)\) satisfies the relation

\[
\widetilde {u (\varphi)} = f * \check {\varphi}
\]

for all \(\varphi\in\mathrm{L}^{p}(\mathrm{G},\mu^{\prime})\) .
% semantic-fragments: legacy-input-seam
\relax{}
